\chapter{相关研究综述与威胁模型}
\label{chap:2}

本章分两部分。第一部分（第 2.1 至 2.5 节）综述相关研究：先给出内存安全攻防方法演进的整体图景（第 2.1 节），然后按内核地址空间布局随机化的三个发展阶段梳理这条路线的现状（第 2.2 节），随后针对持续随机化必须回答的三个核心问题分别综述已有解法及其不足（第 2.3 节），最后介绍本文两项关键技术手段的研究现状，即只执行内存（第 2.4 节）与硬件辅助的控制流完整性（第 2.5 节）。第二部分（第 2.6 至 2.7 节）在综述的基础上给出本文统一的威胁模型，并由威胁模型导出三条贯穿全文的安全不变式。

后续第 3、4、5 章的安全性分析均引用本章第 2.6 节定义的威胁模型与第 2.7 节定义的不变式。

\section{内存安全漏洞攻防方法的演进}

第 1.1 节已经给出内存安全的攻击模型，见图~\ref{fig:1-1}\citep{szekeres2013sok}。本节沿这一模型展开攻防演进的细节，为后续各节提供共同的背景。

最早的内存破坏攻击直接向进程内存中写入可执行的机器码。攻击者利用栈缓冲区溢出覆写返回地址，使其指向注入到缓冲区中的 shellcode，函数返回时即执行攻击载荷。此类攻击对应图~\ref{fig:1-1} 中“篡改代码 → 执行注入的代码”这条路径。

防御方的回应是数据执行阻止（DEP）\citep{paxteam2003pax}：借助页表中的不可执行位（x86 的 NX、ARM 的 XN），使数据页不可执行、代码页不可写，任一内存页不同时具有可写与可执行权限。在内核态，与之对应的机制是内核代码段只读、SMEP（Supervisor Mode Execution Prevention）与 ARM 的 PXN（Privileged eXecute Never）。后两者阻止内核态执行用户态页面上的代码，封堵了“把 shellcode 放在用户态内存、诱使内核跳转过去”这一内核特有的注入路径。

如第 1.1 节所指出的，这些机制在内核中的性质与在用户态不同：它们由内核自己建立和维护，攻击者获得内核态任意写能力之后同样可以关闭。因此内核态的 DEP 类机制提高了攻击成本，却不构成攻击者无法逾越的边界。

代码注入被封死之后，攻击者转向复用程序自身已有的指令片段。Shacham 提出的返回导向编程（ROP）\citep{shacham2007geometry}以 \lstinline[breaklines=false]@ret@ 结尾的短指令序列（gadget）为基本单元，通过在栈上布置一系列 gadget 地址将它们串联执行，并证明在足够大的代码库中这一构造具有图灵完备的表达能力。其后的变体沿两个方向扩展：

\textbf{改变串联手段。} 跳转导向编程（JOP）\citep{bletsch2011jop}以间接跳转指令结尾的片段替代 \lstinline[breaklines=false]@ret@，摆脱了对栈的依赖；面向调用的编程（COOP）\citep{schuster2015coop}以虚函数调用为串联手段，使攻击载荷在控制流图上看起来仍然合法，从而绕过粗粒度的控制流完整性检查。

\textbf{改变载荷构造时机。} 传统 ROP 要求攻击者在发起攻击前就知道 gadget 的地址；即时 ROP（JIT-ROP）\citep{snow2013just}把 gadget 的发现推迟到运行时，详见第 2.1.4 节。

在内核场景下，代码重用攻击的可用素材远比用户态丰富。本文对内核镜像做了一次 gadget 普查作为基线，结果如表~\ref{tab:2-1} 所示。

\begin{table}[!htbp]
    \bicaption{\enspace Linux 内核镜像中可利用 gadget 的分布}{\enspace Distribution of Usable Gadgets in the Linux Kernel Image}
    \label{tab:2-1}
    \centering
    \footnotesize
    \setlength{\tabcolsep}{3pt}
    \renewcommand{\arraystretch}{1.25}
    \begin{tabular}{llrr}
        \hline
        类别 & 终止指令 & 数量 & 占镜像中全部候选序列\\
        \hline
        ROP & \lstinline[breaklines=false]@ret@ & 30 148 & 2.6\%\\
        ROP & \lstinline[breaklines=false]@ret imm@ & 77 067 & 6.5\%\\
        JOP & \lstinline[breaklines=false]@jmp reg@ & 3 015 & 0.3\%\\
        JOP & \lstinline[breaklines=false]@jmp [mem]@ & 18 009 & 1.5\%\\
        COP & \lstinline[breaklines=false]@call reg@ & 3 515 & 0.3\%\\
        COP & \lstinline[breaklines=false]@call [mem]@ & 15 825 & 1.3\%\\
        SYS & \lstinline[breaklines=false]@syscall@ / \lstinline[breaklines=false]@int@ & 697 & 0.1\%\\
        可利用小计 &  & 148 276 & 12.6\%\\
        \hline
    \end{tabular}
\end{table}

按大类归并为经典 ROP 107 215 条、JOP 21 024 条、COP 19 340 条、系统调用类 697 条。其中价值最高的是 \lstinline[breaklines=false]@pop <reg> ; ret@ 一族，它把栈上由攻击者控制的值直接弹入寄存器，是构造攻击链时控制参数的核心原语；镜像中 \lstinline[breaklines=false]@pop rdi ; ret@（控制第一个参数）有 208 条。

这些 gadget 的数量在部署本文机制之后并不会减少，改变的是它们的地址在运行时是否稳定。本文的安全性评估据此设计，比较的是地址的有效期而非 gadget 的计数。

针对代码重用攻击的低开销缓解手段是地址空间布局随机化（ASLR）\citep{paxteam2003pax,shacham2004effectiveness}：在程序装载时随机选择各段的基址，使攻击者无法预知 gadget 的绝对地址。ASLR 的随机熵取决于可用的地址范围与对齐约束，其随机化粒度是整段，段内各函数的相对偏移保持不变。因此攻击者只要泄露段内任意一个地址，就能推算出该段内所有 gadget 的位置。

为提高随机熵并消除这一弱点，研究者提出了一系列细粒度的装载时随机化方法，按粒度可分为四类：内存页粒度\citep{backes2014oxymoron}、函数粒度\citep{kil2006aslp,giuffrida2012enhanced}、基本块粒度\citep{wartell2012stirring}与指令粒度\citep{pappas2012smashing,hiser2012ilr}。自动化软件多样化这一整体思路的设计空间与取舍已有系统梳理\citep{larsen2014sok}，其中面向真实产品的装载时实现以 selfrando 为代表\citep{conti2016selfrando}。粒度越细，随机熵越高，攻击者从单个泄露地址推断其他地址的能力越弱；代价则是重定位信息的规模、装载时间与代码局部性的损失；后一项代价可由编译器在构建期保留必要的重定位信息来摊薄，使随机化在分发或装载时完成而不必在运行期重做静态分析\citep{koo2018compiler}。

这一类方法有一个共同前提：随机化只发生一次。因此它们的安全性完全依赖布局在整个运行期内不被泄露。

信息泄露攻击直接针对上述前提。按泄露手段可分为四类，本文在第 2.6.2 节将它们整理为四类攻击向量。

\textbf{直接读取代码。} JIT-ROP\citep{snow2013just}是这一类的代表：攻击者利用一次内存泄露获得任一代码指针，以该指针所在页为起点递归扫描代码页、反汇编并收集 gadget，最终在运行时构造出攻击载荷。它不需要预先知道完整布局，一个指针即可启动。

\textbf{通过执行结果推测。} 攻击者向猜测的地址发起控制流转移，根据是否崩溃、以何种方式崩溃来判断该地址是否有效。此类方法的关键在于避免崩溃暴露攻击行为。用户态的典型手段有两种：Blind ROP\citep{bittau2014blind}利用网络服务器工作进程崩溃后由守护进程重启且内存布局相同的特性完成持续试探；CROP\citep{kollenda2017crop}复用被攻击程序自身注册的异常处理函数来吸收非法访问异常。

\textbf{这两种手段在内核场景下都不成立。} Blind ROP 依赖 \lstinline[breaklines=false]@fork()@ 语义下父子进程内存布局相同这一前提，内核不存在与之对应的机制；CROP 依赖被攻击程序自行注册的用户态异常处理函数，内核的异常处理路径不由攻击者控制。因此在内核中，攻击者的执行试探会直接触发内核异常，其后果是崩溃或者被内核的异常处理路径捕获，二者都是可观测事件。这一差异对本文有利，是第 4 章检测机制得以成立的一个前提，在第 2.6.2 节定义攻击向量 V2 时会再次说明。

\textbf{通过内存管理接口间接推断。} 攻击者反复申请与释放不同大小的内存块，根据分配是否成功推断虚存空间中未分配区域的分布，进而反推已分配区域的位置\citep{oikonomopoulos2016poking}。此类方法不产生任何非法访存，隐蔽性最强。

\textbf{通过微架构侧信道推测。} 利用缓存访问时序、TLB 状态、预取指令的行为差异等推断地址是否已映射\citep{gruss2016prefetch}。这条路线上的代表性工作包括基于地址翻译缓存时序的早期攻击\citep{hund2013practical}、以事务内存抑制异常的 DrK\citep{jang2016drk}、利用分支目标缓冲冲突的方法\citep{evtyushkin2016jump}，以及在页表隔离部署之后仍然有效的预取侧信道\citep{liu2023entrybleed,lipp2022amd}与针对内核数据结构的侧信道\citep{maar2025kernelsnitch}；瞬态执行攻击为这类泄露提供了更多来源\citep{canella2019transient}。从随机化一侧抵抗侧信道的尝试亦有先例\citep{gens2017lazarus}。近年针对内核地址空间的侧信道工作尤其活跃，例如利用带标签 TLB 打破内核地址空间隔离的 TagBleed\citep{koschel2020tagbleed}，以及专门讨论 ASLR 在微架构攻击下如何加固的 Oreo\citep{song2025oreo}。此类方法同样不产生异常，本文将其列为排除项（第 2.6.4 节）。

上述第一类中的 JIT-ROP 最能说明信息泄露为何足以击穿一次性随机化，其攻击流程如图~\ref{fig:2-1} 所示。攻击者先借内存泄露漏洞得到任意一个代码指针，以它所在的代码页为起点，读取该页内容并反汇编；页内的直接跳转与调用指令给出了指向其他代码页的边指针，攻击者顺着这些边指针递归地读取、反汇编更多代码页，逐步搜集到足够的 gadget；最后用泄露到的真实地址把这些 gadget 拼接成攻击载荷。整个过程发生在程序运行时，一次性随机化在装载时确定的布局对它完全透明。

\begin{figure}[!htbp]
    \centering
    \includegraphics[width=\textwidth,height=0.46\textheight,keepaspectratio]{schematic/ch2-jitrop-disclosure}
    \bicaption{\enspace JIT-ROP 信息泄露攻击流程}{\enspace Workflow of the JIT-ROP Information-Disclosure Attack}
    \label{fig:2-1}
\end{figure}

与随机化并行的另一条防御路线是控制流完整性（CFI）\citep{abadi2005cfi}：静态确定每个间接控制流转移的合法目标集合，在运行时验证实际目标是否属于该集合。CFI 的安全强度取决于目标集合的精度，而精度受限于静态分析在 C 语言上的能力。Carlini 等人\citep{carlini2015bending}通过控制流弯折攻击证明，即便目标集合恰好等于静态可达集合，攻击者仍可沿合法的控制流边构造出恶意的执行序列。

更进一步，数据导向编程（DOP）\citep{hu2016dop}完全不改变控制流，仅通过篡改数据实现图灵完备的计算，从原理上绕过所有 CFI 类防御。这类攻击对应图~\ref{fig:1-1} 中“篡改数据变量 → 纯数据攻击”这条路径（图中已灰显），明确不在本文的防御范围内，理由已在第 1.1 节说明，界定见第 2.6.4 节。

从上述演进可以看到一条清晰的线索：每一项防御机制在某个阶段封死一条路径，攻击者随之转向仍然通畅的路径。本文的工作位于图~\ref{fig:1-1} 中着色的那几条路径，即篡改代码指针、读取代码与代码索引，以及二者汇合而成的控制流劫持路径。所采取的路线是随机化，并以硬件 CFI 补齐随机化在结构上无法覆盖的部分。随机化这一路线的工程化程度与部署障碍已有系统性的梳理\citep{larsen2020adoption}，面向内存泄露的运行期重随机化亦有综述可循\citep{shen2019survey}。

\section{内核地址空间布局随机化相关工作}

将随机化引入内核的研究可以划分为三个阶段。三个阶段并非简单的时间先后，每一阶段都在弥补前一阶段被攻破的假设。第一阶段是一次性随机化，它假定代码布局在运行期不被攻击者获知，信息泄露与侧信道攻破了这一假设；第二阶段是持续随机化，它以周期性重排使已泄露的地址迅速失效，但全量重排的代价使其难以覆盖整个内核；第三阶段是部分随机化与触发式随机化，它通过限制保护范围或触发时机把代价压到可接受的水平，代价是留下布局固定的非随机化区域。以下三小节依次讨论这三个阶段。

\subsection{一次性随机化}

面向内核的装载时随机化最早由研究工作提出\citep{stanley2013improved}；Linux 自 3.14 版本起在主线内核中提供 KASLR\citep{linux2024kernel}，在系统启动时随机选择内核映像的装载基址。这一机制的随机熵受两方面限制：可用于放置内核映像的虚拟地址范围有限，以及内核映像必须按大页边界对齐。二者共同决定了可选位置的数量。

除信息泄露之外，内核映像的基址还可由微架构侧信道恢复，这一方向的工作持续推进\citep{gruss2017kaiser,canella2020kaslr}；即便完全没有信息泄露，攻击者也可构造位置无关的代码重用载荷\citep{goktas2018pir}。这一阶段的共同弱点是随机化只发生一次。第 2.1.4 节的任一类信息泄露手段一旦成功，防御即告失效，且不再恢复。

\subsection{持续随机化}

针对上述弱点，研究者提出在运行期周期性地重新排列代码布局。这一方向的工作主要在用户态展开，代表性系统包括 TASR\citep{bigelow2015tasr}、Shuffler\citep{williams2016shuffler}、Remix\citep{chen2016remix}、RuntimeASLR\citep{lu2016runtimeaslr}、CodeArmor\citep{chen2017codearmor}、ReRanz\citep{wang2017reranz}与 MARDU\citep{jelesnianski2020mardu}。另一条并行的思路不改变代码位置而随机化执行路径本身，以克隆函数并在每个调用点随机选择副本的方式挫败按地址预测的攻击链\citep{davi2015isomeron}。此外还有 Mixr\citep{hawkins2017mixr} 等以二进制为对象的工作。它们的差异集中在三点，即何时触发随机化、如何更新代码索引、以及如何处理正在执行的控制流，将在第 2.3 节展开比较。

上述方案对代码索引的处理都依赖于把程序改造为可重排的形式，这一能力在无源码场景下由静态二进制重写工具提供，代表性工作以低开销的直接改写\citep{hawkins2017zipr}与完全布局无关的重编译\citep{williamsking2020egalito}为两条路线。Shuffler 的实验给出了这一阶段最重要的一个量化结论：其 50 ms 的洗牌周期远短于已知 JIT-ROP 攻击所需的 2.3 至 378 秒。Ahmed 等人随后的系统测量\citep{ahmed2020methodologies}把这个经验判断精确化为可复现的方法，通过测量 gadget 集合随扫描时间的收敛速度，给出随机化间隔的上界为 1.5 至 3.5 秒。本文第 4 章不采用这一以扫描时间反推间隔上界的口径，因为触发式随机化的时机由检测事件决定，而不由间隔参数决定。

这一阶段的困难在于代价。如第 1.2 节所述，代码索引更新与执行流暂停两项开销在内核中被显著放大，这直接导致了第三阶段的出现。

\subsection{部分随机化与触发式随机化}

为使持续随机化在内核中可行，近年的工作沿两个方向缩小代价。

\textbf{缩小随机化的范围。} Adelie\citep{nikolaev2022adelie,nadeem2020kernelmodule}只对可加载内核模块与设备驱动应用持续随机化。其技术路线包括：将模块编译为位置无关代码并为每个模块维护只读的全局偏移表；以 \lstinline[breaklines=false]@schedule_on_each_cpu()@ 配合 RCU 宽限期确保重随机化期间无执行流停留在被移动的代码中；以 20 ms 的固定周期由专用内核线程驱动重随机化。其安全分析给出的攻击成功概率为 2$^{-44}$（考虑 4 KB 页对齐之后），并与 Shuffler、CodeArmor 使用 32 位偏移所对应的 2$^{-19}$ 作对比。

\textbf{缩小随机化的时机与对象。} 另一条路线是放弃固定周期，只在观测到实际的泄露行为时才随机化，并把随机化的对象收窄到被泄露的那一部分代码。PtrProxy\citep{luo2025ptrproxy}在 AArch64 上采取这一思路，以代码指针的实际泄露作为触发条件，并用跳板承载指针代理，跳板本身的泄露亦触发检测。这条路线的价值在于把代价与攻击行为绑定，在无攻击的时段几乎不付出开销；其难点则在于泄露行为的观测手段本身要足够廉价且足够完备。\textbf{目前这一路线的工作均以用户态进程为保护对象，尚无面向内核 built-in 代码的实现。}

表~\ref{tab:2-3} 汇总了这两个方向上的代表性工作，列的划分对应第 1.3 节提出的三个核心问题。

\begin{table}[!htbp]
    \bicaption{\enspace 持续随机化代表性工作的横向对比}{\enspace Comparison of Representative Continuous Re-randomization Systems}
    \label{tab:2-3}
    \centering
    \scriptsize
    \setlength{\tabcolsep}{3pt}
    \renewcommand{\arraystretch}{1.25}
    \hyphenpenalty=10000\exhyphenpenalty=10000\emergencystretch=1em
    \begin{tabular}{>{\raggedright\arraybackslash}p{\dimexpr 0.1440\textwidth-2\tabcolsep\relax}>{\raggedright\arraybackslash}p{\dimexpr 0.0720\textwidth-2\tabcolsep\relax}>{\raggedright\arraybackslash}p{\dimexpr 0.0871\textwidth-2\tabcolsep\relax}>{\raggedright\arraybackslash}p{\dimexpr 0.0758\textwidth-2\tabcolsep\relax}>{\raggedright\arraybackslash}p{\dimexpr 0.1061\textwidth-2\tabcolsep\relax}>{\raggedright\arraybackslash}p{\dimexpr 0.1970\textwidth-2\tabcolsep\relax}>{\raggedright\arraybackslash}p{\dimexpr 0.1061\textwidth-2\tabcolsep\relax}>{\raggedright\arraybackslash}p{\dimexpr 0.0947\textwidth-2\tabcolsep\relax}>{\raggedleft\arraybackslash}p{\dimexpr 0.1023\textwidth-2\tabcolsep\relax}}
        \hline
        系统 & 年份 · 会议 & 保护对象 & 粒度 & 何时随机化 & 如何更新代码索引 & 执行流处理 & 非随机化区域 & 报告开销\\
        \hline
        TASR & 15 · CCS & 用户态进程 & 函数 & 成对 I/O 系统调用之间 & 编译期标记，运行时更新全部代码指针 & 系统调用边界 & 无 & 2.1\%\\
        Remix & 16 · CODASPY & 用户态进程 & 函数内基本块 & 固定间隔 & 函数位置不变，无需更新 & 无需暂停 & 无 & 2.8\%\\
        RuntimeASLR & 16 · NDSS & 用户态进程 & 模块 & \lstinline[breaklines=false]@fork()@ 时 & 动态污点分析追踪指针 & 仅在 fork 点 & 无 & 记录阶段数百倍\\
        Shuffler & 16 · OSDI & 用户态进程 & 函数 & 固定 50 ms & 代码指针表下标 & 保留旧副本，异步洗牌 & 无 & 14.9\%\\
        CodeArmor & 17 · EuroS\&P & 用户态进程 & 整段 & 系统调用 & 地址空间抽象层 & 抽象层转换 & 无 & 6.9\%\\
        ReRanz & 17 · VEE & 用户态进程 & 基本块 & I/O 字节累积量 & 进程级虚拟机跳板 & 虚拟机接管 & 无 & 6\%\\
        SafeHidden\citep{wang2019safehidden} & 19 · USENIX Sec & 用户态安全区域 & 区域 & 检测到试探 & 不适用 & 不适用 & 陷阱区域 & 2.75\%\\
        MARDU & 20 · SYSTOR & 用户态进程 & 函数 & 泄露时 & 跳板 & 保留旧副本 & 无 & 5.5\%\\
        \textbf{Adelie\citep{nikolaev2022adelie}} & \textbf{22 · ASPLOS} & \textbf{内核模块与驱动} & \textbf{模块} & \textbf{固定 20 ms} & \textbf{PIC + 每模块只读 GOT} & \textbf{调度器 + RCU 宽限期} & \textbf{只读 GOT 约束跨模块指针} & \textbf{< 2\%}\\
        Dbox\citep{li2024dbox} & 24 · CCS & 不可信驱动 & 整个驱动的虚拟空间 & 捕获到进入驱动的可疑控制流 & 不搬移物理代码，只改虚拟地址映射，因而无需更新索引 & 无需暂停 & 虚拟盒隔离 & < 3\%（CPU 负载 40\% 以下）\\
        PtrProxy & 25 · China Comm. & 用户态进程（AArch64） & 基本块 & 代码指针实际泄露 & 跳板承载指针代理 & 指针代理转发，无可感知暂停 & 跳板泄露触发检测 & 1.52\%\\
        \textbf{本文} &  & \textbf{内核 built-in 代码} & \textbf{函数} & \textbf{检测到代码探测} & \textbf{函数级 PIC + 固定跳板} & \textbf{跳板计数精确暂停，无副本} & \textbf{随机化辅助的指针认证} & \textbf{见第 3、4、5 章}\\
        \hline
    \end{tabular}
\end{table}

这一阶段的代价是范围。部分随机化带来两个后果：其一，built-in 内核代码完全不受保护，而第 1.2 节已经说明这部分代码既无法被这类方案覆盖，又承载了两成以上的内核 CVE；其二，非随机化区域的布局对攻击者而言是已知且固定的，本身成为新的攻击面（详见第 2.3.4 节）。

本文的出发点即在于此：把持续随机化扩展到 built-in 内核代码，同时处理由此产生的非随机化区域问题。

\section{持续随机化的三个核心问题}

如第 1.3 节所述，持续随机化机制必须回答三个问题。本节按这三个问题组织已有工作，每一小节的结构相同：已有解法的分类、各自的不足、本文对应章节的出发点。

\subsection{代码索引更新}

代码位置改变后，所有指向代码的引用都需要相应更新。已有方法按更新时机可分为三类。

\textbf{在赋值时更新，即指针追踪类。} 这类方法在运行时追踪所有代码指针的传播，随机化时逐一改写其值。TASR\citep{bigelow2015tasr}修改编译器收集代码指针的初始位置并动态监控关键数据结构的传播；OS-level ASR\citep{giuffrida2012enhanced}在编译器中间表示上插桩以跟踪指针传播；RuntimeASLR\citep{lu2016runtimeaslr}采用动态污点分析追踪指针。这类方法有两处困难：需要源代码，因而无法支持闭源与遗留二进制；指针别名与类型转换使追踪难以完备，且动态追踪的开销极高，RuntimeASLR 在记录阶段引入的开销达数百倍。

\textbf{在赋值时替换语义，即间接层类。} 这类方法从一开始就把代码指针替换为指向某个间接结构的索引。Oxymoron\citep{backes2014oxymoron}将函数指针替换为指向受保护的类 GOT 区域的跳板；Shuffler\citep{williams2016shuffler}将代码指针替换为全局代码指针表的下标；CodeArmor\citep{chen2017codearmor}引入一层地址空间抽象；MARDU\citep{jelesnianski2020mardu}以跳板作为随机化的间接层。这类方法把索引更新量从“所有指针”降至“所有跳板”，代价是改变了代码指针的语义：指针的值不再是地址，因此不支持自引用代码。栈展开库通过返回地址读取其指向的指令编码以判断特殊指令模式，这类用法在语义被替换后即告失效。

\textbf{在使用时翻译，即动态翻译类。} 这类方法不改变指针的值，只在指针被解引用发生控制流转移时进行地址翻译。ReRanz\citep{wang2017reranz}以进程级虚拟机实现这一思路，从而支持闭源二进制。其代价是需要一层运行时翻译，且翻译表的规模随间接分支目标的多样性增长。

\textbf{回避问题的一类。} Remix 只重排函数内部的基本块而保持函数本身的位置不变，因此代码索引根本不需要更新。这一取舍换来了极低的开销，代价是函数入口地址仍然固定：攻击者若已知函数地址，仍可将整个函数当作一个大 gadget 使用，函数粒度以上的布局信息也不会因随机化而失效。

\textbf{内核场景。} Adelie 采用间接层类思路的变体：模块编译为位置无关代码消除模块内部的索引，模块之间的引用统一经由只读全局偏移表，随机化时只需更新表项。这一方案的可行性依赖一个前提，即模块与核心内核之间存在有限且明确的接口。built-in 代码不具备这一前提，核心子系统之间的调用关系复杂且动态，无法用一张表把跨区域引用全部收拢。

\textbf{不足与本文的出发点。} 三类方法在 built-in 内核代码上都遇到障碍：指针追踪类的开销与完备性问题在千万行规模的内核上被放大；间接层类因改变指针语义而与内核中大量将函数地址当作数值使用的代码模式冲突；动态翻译类的翻译层开销在内核热路径上不可接受。回避问题的一类则以函数位置固定为代价，与本文“使已泄露的函数地址失效”这一目标直接冲突。本文第 3 章的出发点是把索引分层处理，函数粒度内用位置无关代码消除，函数粒度间用固定位置的跳板集中，并把“数值还是指针”的判定推迟到解引用点，从而在不改变指针值的前提下把一次随机化需要改写的代码索引由与调用点数量同阶降到与函数数量同阶。

\subsection{执行流的安全暂停}

在改变代码位置期间，必须确保没有执行流正停留在被移动的代码中，否则将跳转到无效地址。已有方法可分为两类。

\textbf{保留旧副本。} 随机化时不回收旧代码，令仍在旧代码中执行的线程自然结束，因而无需暂停任何线程。Shuffler 与 MARDU 采用这一路线。其代价是旧副本在被回收前仍是攻击者可用的 gadget 来源。这与本文的防御目标直接冲突：本文的整个论证建立在“随机化使已泄露的信息过期”之上，而旧副本恰恰使信息不过期。

\textbf{暂停执行流。} 随机化前先使所有可能受影响的执行流停下。TASR 选择在系统调用边界进行随机化，此时线程状态明确；Adelie 以 \lstinline[breaklines=false]@schedule_on_each_cpu()@ 配合 RCU 宽限期，确保重随机化期间无执行流停留在被移动的模块代码中。

\textbf{不足与本文的出发点。} 保留副本的方案以安全换取免暂停。暂停类方案在内核中的困难是非抢占上下文：中断处理程序不接受调度；RCU 宽限期只能确保读端临界区结束，无法覆盖中断处理程序；强制暂停持有自旋锁的代码会导致其他处理器死锁等待该锁。本文第 3 章的出发点是利用跳板函数作为随机化区域边界上的观测点，通过入口与出口计数精确得到当前位于随机化区域内的执行流，从而只暂停必要的执行流，并对非抢占上下文采用推迟随机化而非强制暂停。

\subsection{随机化的触发策略}

触发策略决定了随机化的频率，直接决定安全性与开销的平衡点。已有策略可分为四类。

\textbf{固定时间间隔。} Shuffler 采用 50 ms，CodeArmor 可低至 55 μs，Adelie 采用 20 ms，Remix 与 OS-level ASR 同样采用固定周期。这类策略实现简单，却面临第 1.3.3 节所述的两难：间隔过短则开销不可承受，过长则留下攻击窗口；且无论间隔取值如何，代价被平摊在整个运行期，而收益只在“泄露已发生、劫持尚未完成”这一极窄窗口内产生。

\textbf{特定事件驱动。} TASR 在成对的输入与输出系统调用之间触发，其依据是任何输出类系统调用都可能被用于泄露、任何输入类系统调用都可能被用于上传载荷；RuntimeASLR 在父进程创建子进程时触发，针对的是基于克隆的试探攻击。这类策略把触发与程序行为关联起来，但关联的是“可能被利用的行为”而非“实际发生的攻击行为”，对 I/O 密集型负载会触发过频。

\textbf{输出量驱动。} ReRanz 提出以数据传输的累积量作为触发条件，只有当应用输出足够多的内容时才进行随机化，从而在 I/O 密集型场景下显著优于 TASR 的策略。这一策略仍然是对“泄露可能已经发生”的间接估计。

\textbf{泄露驱动。} 一类工作在观测到代码被读取时随机化被读到的那一部分\citep{morton2017zombie}，或在检测到地址探测后收回被探测代码的执行权限\citep{li2022mprobe}、重排被探测函数内部的代码块\citep{li2024randomize}。这一路线的工作均以用户态进程为保护对象，第 4.2.2 节对其与本文的差别有进一步分析。

\textbf{探测行为驱动。} SafeHidden 以检测到的内存试探行为触发安全区域的重定位；PtrProxy 以代码指针的实际泄露作为触发条件。\textbf{这一类策略与本文第 4 章的思路同属一条路线}，其差异见下。

\textbf{不足与本文的出发点。} 前三类策略的共同点是都没有利用“代码是否真的被探测过”这一信息，因而只能在“宁可多随机”与“冒险少随机”之间取舍。第四类策略已经利用了这一信息，但其保护对象均为用户态进程或用户态的安全区域，尚无工作把它用于内核自身的 built-in 代码。本文第 4 章的出发点是把这一思路带入内核态，并解决由此产生的两个内核特有问题：内核代码页存在多重映射，使基于页表权限的实现可被绕过；内核中存在合法的自读代码路径，纯防护式的只执行内存必须为其逐一豁免。

\subsection{非随机化区域的处理}

部分随机化必然留下非随机化区域，其布局对攻击者是已知且固定的。已有工作对此的处理相对薄弱。

Adelie 的处理方式是限制两类指针的可利用性：指向模块内部的指针在模块移动后自动失效；指向非随机化内核代码的指针被约束为只能经由只读的全局偏移表访问，且该表以 RIP 相对寻址访问、随模块一同移动，因而无法被伪造。这一方案依赖模块与内核之间的接口有限且明确。

更一般地，非随机化区域的保护属于控制流完整性的范畴，相关技术在第 2.5 节综述。

\textbf{不足与本文的出发点。} Adelie 的方案不适用于 built-in 场景，built-in 代码之间不存在有限的接口集合，无法用一张只读表把跨区域引用全部收拢。通用的 CFI 方案在内核中则面临第 2.5 节所述的精度与开销约束。本文第 5 章的出发点是利用“随机化区域内的控制流因目标地址不可预测而天然不可劫持”这一性质，只对跨区域的间接转移施加硬件 CFI 验证，并借随机化打破指针认证的上下文传播链。

\section{只执行内存技术}

只执行内存（Execute-Only Memory，XOM）使代码页可执行但不可读，从而在图~\ref{fig:1-1} 的第 2 步上直接封堵“用指针读代码”这条路径。本节综述其实现路径与在内核中的困难。图~\ref{fig:2-2} 对照三条实现路径，其中页表权限模拟作为负面例子列出，第 4 章的两条实现路径取自后两条。

\begin{figure}[!htbp]
    \centering
    \includegraphics[width=\textwidth,height=0.46\textheight,keepaspectratio]{schematic/ch2-xom-three-paths}
    \bicaption{\enspace 只执行内存的三条实现路径}{\enspace Three Implementation Paths for Execute-Only Memory}
    \label{fig:2-2}
\end{figure}

\subsection{实现路径}

\textbf{软件模拟。} 在硬件不提供只执行权限的平台上，可通过页表权限的动态切换或分离的指令与数据 TLB 来模拟。XnR\citep{backes2014xnr}始终只保持少量代码页可读，其余置为不可访问，通过缺页异常在滑动窗口内换入换出；HideM\citep{gionta2015hidem}利用分离 TLB 使取指与取数看到不同的物理页。在不提供只执行权限的嵌入式平台上，另有以软件掩码约束读操作\citep{braden2016leakage}、以非特权访存指令配合内存保护单元\citep{kwon2019uxom}、以及以调试监视比较器实现只执行\citep{shen2020picoxom}的做法；把只执行改造到没有源码的二进制上还须先把嵌在代码中的数据搬离\citep{chen2017norax}。这类方法的开销取决于窗口大小与异常频率。更根本的问题在于多重映射，见第 2.4.2 节。

\textbf{二级地址翻译。} Intel 的扩展页表（EPT）\citep{intel2023sdm}为二级地址翻译提供了独立的权限位，可以配置为可执行但不可读，且该权限由虚拟机监控器维护、客户机内核无法修改。Readactor\citep{crane2015readactor}即基于 EPT 实现用户态的只执行内存，并配合代码指针隐藏抵御 JIT-ROP。这条路径的优点是作用于物理页面，天然覆盖同一物理页的全部虚拟映射；代价是要求系统运行在虚拟化环境中。

\textbf{调试观察点。} AArch64 的调试观察点提供了另一条与页表无关的路径。观察点只对数据访问触发，不对取指触发，这恰好就是只执行的语义，即允许执行、捕获读取。其地址掩码域可掩蔽 3 至 31 位地址，因此单个观察点最大可覆盖 2 GB 的自然对齐区域，只要把受保护的代码集中放置在一段连续且按 2 的幂对齐的地址空间内，观察点的数量就不构成覆盖上的限制。这条路径不要求虚拟化层，且与二级地址翻译一样，其配置位于调试寄存器中，独立于内核所维护的页表之外（第 1.1 节所述内核第一条约束下的另一条路径）。

\textbf{检测而非阻断。} 这一类方法允许代码读取发生，但使读到的内容失去价值。Heisenbyte\citep{tang2015heisenbyte}提出破坏性代码读取，借助 EPT 捕获对代码页的读取，向攻击者返回原始内容的同时破坏实际执行的那一份，使已泄露的代码无法再被执行；NEAR\citep{werner2016near}采取类似思路。内核场景下针对即时代码重用的系统性防护亦有工作\citep{pomonis2019krx}。第 4 章属于这一路线，其响应方式与已有工作不同：已有工作在检测到读取后破坏被读代码或禁止其后续执行，因而必须为每一处被读代码维护副本或降级其语义；第 4 章的响应是触发一次针对性的重随机化，不改变被读代码的任何语义。二者的差异见第 4.2 节。

\subsection{内核场景下的两个困难}

\textbf{代码页的别名映射。} 内核代码在物理内存中通常同时存在于内核映像的线性映射与直接映射之下。以页表权限实现的只执行内存若只作用于其中一份映射，攻击者可经由另一份映射绕过，且这一绕过不需要篡改任何权限。基于二级地址翻译的实现因作用于物理页面而天然覆盖所有映射，但它要求存在虚拟化层。本文在 x86 上改用管理态保护键：保护键记录在页表项中，只要让同一物理页的两份映射使用同一个键，别名映射即被一并覆盖，无须虚拟化层，详见第 4.4.1 节。ARM 上的观察点作用于虚拟地址区间，因此需要额外确认监控区间是否覆盖了同一物理页的其他虚拟映射，这一点在第 4.4.1 节说明。

\textbf{合法的自读代码路径。} 内核中存在若干正当地读取自身代码的机制，纯防护式的只执行内存必须为每一条路径逐一豁免，而豁免清单既难以完备，其本身也构成攻击面。这些路径包括：\lstinline[breaklines=false]@kprobes@ 在插桩时读取并解析原指令；\lstinline[breaklines=false]@ftrace@ 在启用或禁用跟踪点时读取并改写调用指令；静态键与 alternative 机制在运行时改写代码；\lstinline[breaklines=false]@livepatch@ 在应用补丁时读取原函数；崩溃路径上的 \lstinline[breaklines=false]@show_opcodes()@ 会转储故障地址附近的指令字节。

\textbf{一个有利的观察是：这些路径全部集中在离散的事件上，而非稳态运行的持续行为。} 内核在正常执行系统调用、处理中断、调度进程时并不读取自身的指令；异常修复表查询读取的是表本身而非代码，栈回溯读取的是 ORC 或帧指针元数据而非指令。因此上述清单中的每一项都对应一个明确的、可被登记的触发时刻：插桩注册、跟踪点开关、静态键翻转、补丁应用、以及崩溃转储。

本文第 4 章的出发点由此形成：既然合法读取路径都对应离散且可识别的事件，就不再试图阻断读取，而是把读取当作信号，合法路径只需被登记为已知来源，无需被逐一豁免。

\section{硬件辅助的控制流完整性}

本节综述硬件辅助 CFI 的三类机制及其精度上限。第 1.1 节已给出结论，本节给出依据。图~\ref{fig:2-3} 按等价类的划分方式对照三类机制，图中的数字为本文在 Linux 6.8 上的实测值，各自的测量配置与口径见以下三小节。

\begin{figure}[!htbp]
    \centering
    \includegraphics[width=\textwidth,height=0.46\textheight,keepaspectratio]{schematic/ch2-cfi-equivalence-class}
    \bicaption{\enspace 三类硬件 CFI 机制的等价类划分}{\enspace Equivalence-Class Partitioning of Three Classes of Hardware CFI}
    \label{fig:2-3}
\end{figure}

\subsection{落地点标记类}

Intel 的控制流强制技术（CET）中的间接分支跟踪（IBT）与 ARM 的分支目标识别（BTI）采用相同思路：以一条特殊指令（\lstinline[breaklines=false]@endbr64@ 或 \lstinline[breaklines=false]@BTI@）标记合法的间接分支目标，硬件在每次间接转移后检查目标处是否存在该标记，否则触发异常\citep{shanbhogue2019cet,arm2023armv8}。

这一机制的精度上限是明确的：所有被标记的位置构成单一的等价类，任何间接调用都可以合法地跳向其中任意一个，是一种已知可被绕过的宽松策略\citep{gaidis2023fineibt}。内核社区为此引入了 \lstinline[breaklines=false]@objtool@ 的封印机制，静态分析出不可能被间接调用的函数，把其入口处的 \lstinline[breaklines=false]@endbr64@ 改写为 \lstinline[breaklines=false]@nop@，从而缩小可用目标集合；收缩之后的集合仍然很大。

在此基础上，FineIBT\citep{gaidis2023fineibt}通过在调用点设置类型哈希、在被调用点检查，把 IBT 提升为细粒度策略。以函数类型划分等价类的策略已在 Android 内核上大规模部署，其公开报告给出了这一策略在真实内核上的精度\citep{tolvanen2018androidcfi}。CET 的另一半是影子栈，它与前向边的 CFI 互补，负责保护返回地址\citep{burow2019shadow}；把代码指针本身放入受隔离区域是同一问题的另一条解法\citep{kuznetsov2014cpi}。

在 Linux 6.8（x86-64，\lstinline[breaklines=false]@x86_64_defconfig@ + \lstinline[breaklines=false]@CONFIG_X86_KERNEL_IBT@）上实测：编译器共插入 13 322 条 \lstinline[breaklines=false]@endbr64@，\lstinline[breaklines=false]@objtool@ 封印其中 4 415 条，剩余 8 907 条，即 IBT 的等价类大小为 8 907。这个数字随内核配置规模变化——在本项目实际使用的完整配置上升至 38 180，因此报数时须附配置。ARM64 的 \lstinline[breaklines=false]@defconfig@ 默认未开启 \lstinline[breaklines=false]@CONFIG_ARM64_BTI_KERNEL@，内核自身几乎不含 BTI 落地点，故 ARM 平台的落地点标记基线不单独给出，CFI 精度对比以类型哈希类为主。

\subsection{签名验证类}

ARM 指针认证（PA）在指针的未使用高位中嵌入一个密码学签名，签名的输入包括指针值、一个 64 位的上下文与一个软件不可读的硬件密钥。签名与验证由专用指令完成，开销接近于零。

PA 的精度取决于上下文的选取。若每个调用点使用互不相同的上下文，则在某调用点签名的指针无法在另一调用点通过验证，等价类大小为一。但这一理想状态无法达到，原因是上下文传播：当多个调用点可能调用同一函数时，这些调用点必须使用相同的上下文，否则被调用方无法验证；该约束沿调用图向外扩散，最终把大量调用点并入同一上下文。Liljestrand 等人测得 iOS 13 中约有 6 500 个调用点共享同一上下文\citep{liljestrand2019pac}。

围绕 PA 的系统性工作包括 PARTS\citep{liljestrand2019pac}、PACStack\citep{liljestrand2021pacstack}与 PACTight\citep{ismail2022pactight}等，面向内核的 PA 应用亦有若干工作。这些工作或改进上下文的构造方式，或扩展保护对象，却都未能从根本上打断传播链，因为传播的根源是调用图本身的结构。

\textbf{攻击面。} PACMAN\citep{ravichandran2022pacman}证明 PA 的验证结果可以经由推测执行侧信道被读出：攻击者把认证指令置于推测执行路径上，通过微架构侧信道观测验证是否成功，从而在不产生崩溃、不留下日志的前提下爆破签名。这一攻击已在商用处理器上跨特权级演示。此外，若程序中存在攻击者可控上下文的签名原语，指针重用与上下文碰撞的分析同样失效。本文第 5 章的安全性论证以“攻击者无法获得 PA 验证预言机”为显式前提，该前提及其脆弱性在第 5.6 节单独说明。

\subsection{类型哈希类}

以函数原型的哈希作为等价类的划分依据，调用点检查目标函数的类型签名是否匹配。Linux 内核中的 kCFI\citep{clang2022kcfi}（\lstinline[breaklines=false]@CONFIG_CFI_CLANG@）即采用这一策略。其精度介于前两类之间。在 Linux 6.8 的 ARM64 镜像上按函数类型签名分组实测：56 470 个有实体的函数落入 17 883 个不同的类型签名，\textbf{平均每个等价类只有 3.16 个成员，但最大的等价类却有 2 139 个成员}。最大的几个等价类均来自最通用的原型，如 \lstinline[breaklines=false]@int (*)(void)@（2 139 个）、\lstinline[breaklines=false]@void (*)(void)@（2 069 个）、\lstinline[breaklines=false]@long (*)(struct device*, struct device_attribute*, char*)@（1 266 个）。\textbf{决定安全性的正是这个 2 139，而非平均值 3.16}\citep{burow2017cfi}：一个类型为 \lstinline[breaklines=false]@int (*)(void)@ 的间接调用，在 kCFI 之下有两千余个合法目标可选。

需要说明测量口径。上述统计基于 DWARF 调试信息中的函数原型，覆盖“所有有实体的函数”，而 kCFI 的等价类严格来说只应包含可能被间接调用的函数，因此这里的分母与最大等价类都偏大；更贴近编译器实际划分的做法是统计 Clang 生成的 \lstinline[breaklines=false]@__kcfi_typeid_@ 符号，但它要求 \lstinline[breaklines=false]@CONFIG_CFI_CLANG@（Clang ≥ 16），本文测量所用工具链未满足，故采用前一种口径并如实标注其偏宽。该数据同时回填第 1.1 节。第 5.5.2 节的对比基线采用更严格的口径，只统计被取过地址的函数并折叠枚举与匿名结构体，在 Linux 6.8 上得到的最大等价类为 141，两个数字的差别全部来自统计口径。类型签名与架构基本无关（同一份 C 源码），故此分布可代表两个平台。

\subsection{小结}

三类机制的精度上限有一个共同根源：硬件只能校验一个可在常数时间内判定的谓词，即有无标记、签名是否匹配、哈希是否相等，而“这个调用点在此刻是否应当跳到这个函数”是上下文相关的，超出了硬件可判定的范围。因此硬件 CFI 的精度上限是机制问题而非实现问题，提升精度必须依靠硬件之外的信息。

本文第 5 章的出发点正在于此：随机化提供了这一额外信息，随机化区域内的目标地址不可预测，因而进入该区域的转移无需验证；把高扇入函数移入随机化区域，即可从中间打断 PA 的上下文传播链。

\section{本文的威胁模型}

本节给出贯穿全文的威胁模型。第 3、4、5 章不再各自定义威胁模型，统一引用本节。

\subsection{攻击者能力}

本文假设攻击者具备如下能力：

\begin{enumerate}
    \item \textbf{任意地址读。} 攻击者可利用内核中的内存安全漏洞读取任意内核虚拟地址的内容。这是第 2.1.4 节各类信息泄露手段的基础。
    \item \textbf{任意地址写。} 攻击者可写入任意内核虚拟地址。这一能力使得内核自身维护的防护位（SMEP、SMAP、代码段只读）不构成可信边界，如第 2.1.1 节所述。
    \item \textbf{触发任意系统调用。} 攻击者可从用户态以任意参数反复调用系统调用，从而驱动内核执行、申请与释放内存、创建线程。
    \item \textbf{已知内核二进制。} 攻击者拥有与目标系统相同版本的内核映像，因而掌握所有代码的相对偏移、gadget 的完整清单与全部符号信息。攻击者唯一不知道的是随机化区域内代码在当前时刻的实际地址。
\end{enumerate}

假设 1 与 2 对应图~\ref{fig:1-1} 中第 1、2 步以及“篡改数据指针”的放大回边。如第 1.1 节所述，本文把它作为攻击者能力的来源而非防御目标。

\subsection{四类攻击向量}

基于上述能力，攻击者可通过以下四类向量获取随机化区域的布局信息或绕过随机化保护。这一分类在第 4、5 章的实验中逐条对应验证。

\textbf{V1 直接读取代码。} 攻击者利用任意读能力直接读取随机化区域的代码内容，反汇编后收集 gadget。JIT-ROP 是其典型形态。

\textbf{V2 通过执行结果推测。} 攻击者向猜测的地址发起控制流转移，或对疑似地址发起访问，通过是否产生异常、产生何种异常来判断猜测是否正确。

\begin{quote}
需要说明该向量在内核中的特殊形态。如第 2.1.4 节所述，用户态的两种隐蔽试探手段在内核中都不成立：Blind ROP 依赖 \lstinline[breaklines=false]@fork()@ 语义，内核无对应机制；CROP 依赖被攻击程序自行注册的异常处理函数，内核的异常处理路径不由攻击者控制。因此内核中的执行试探必然产生内核可见的异常，其后果或是系统崩溃，或是被内核异常路径捕获。\textbf{这使 V2 在内核中比在用户态更易被观测，是第 4 章检测机制的一个有利前提。} 攻击者可能的规避方式是把试探限制在极低的频率上，以混入内核正常运行中偶发的异常噪声。第 4 章的机制对此不依赖频率判据：随机化区域的当前地址只存在于跳转桩的分支立即数中，攻击者要拿到可用的地址就必须读取跳转桩区，而该区域处于只执行保护之下，任何一次读取都产生硬件事件并立即触发随机化。试探放得再慢，每一次读取仍然各自触发一次，探得的地址在下一次使用之前即已失效。
\end{quote}

\textbf{V3 通过代码索引推断。} 攻击者不直接读取代码，而是读取内存中保存的指向代码的索引，包括栈上的返回地址、结构体中的函数指针、跳转表项，并由此推断代码位置。

\textbf{V4 利用非随机化区域。} 攻击者放弃获取随机化区域的布局，转而完全在非随机化区域内构造攻击链。此向量又分两种形态：V4-A 在非随机化区域内构造完整的 ROP 或 JOP 链；V4-B 劫持跨区域的间接转移，借合法的跨区域路径进入随机化区域内的目标。

四类向量的应对方式与三章工作的对应关系见图~\ref{fig:2-4}。图中实心方块表示该向量在结构上被消除，空心圆表示以运行时的检测或验证覆盖，虚线表示某一章的覆盖依赖于另一章提供的机制。图中所列的三条安全不变式在第 2.7 节给出。

\begin{figure}[!htbp]
    \centering
    \includegraphics[width=\textwidth,height=0.46\textheight,keepaspectratio]{schematic/ch2-vectors-and-invariants}
    \bicaption{\enspace 四类攻击向量、三条安全不变式与三章工作的对应关系}{\enspace Four Attack Vectors, Three Security Invariants, and Their Correspondence to Chapters 3 to 5}
    \label{fig:2-4}
\end{figure}

\subsection{信任基础}

本文的防御建立在以下可信前提之上：

\begin{enumerate}
    \item \textbf{安全启动链可信。} 系统启动至内核随机化机制初始化完成之前，攻击者无法介入。因此初始的随机化布局、跳板函数的初始内容与硬件密钥的初始化均可信。
    \item \textbf{处理器按架构手册规定的语义工作。} 管理态保护键的权限位、硬件调试观察点、指针认证指令的行为符合规范，不存在使其失效的未公开勘误。
    \item \textbf{攻击者在劫持控制流之前不能任意改写内存权限。} 本文的两条只执行路径都由内核自身维护，管理态保护键的权限状态保存在每处理器核的寄存器中，观察点的配置保存在调试寄存器中，二者都不经普通页表；该假设保证它们在被绕过之前不会先被关闭。
    \item \textbf{指针认证的硬件密钥不可读。} 密钥由硬件维护，软件无法读出。
\end{enumerate}

\subsection{明确的排除项}

以下内容不在本文的防御范围内。

\textbf{纯数据攻击。} 篡改数据而不改变控制流的攻击（例如直接改写进程凭据结构中的 uid 字段、数据导向编程）与代码布局无关：这类攻击既不需要知道任何 gadget 的位置，也不发生非法的控制流转移，因此以“使代码位置不可预测”为原理的防御在结构上对其无效。这是随机化路线的固有边界，而非本文方案的不足。防御此类攻击需要数据流完整性或内存安全强制等正交手段，可与本文机制叠加部署。

\textbf{微架构侧信道。} 利用缓存时序、TLB 状态、预取行为或推测执行推断地址的攻击不在防御范围内。这类攻击不产生本文检测机制可观测的事件。特别地，这意味着第 2.5.2 节所述的 PACMAN 类攻击不被本文阻止，第 5 章的安全性论证以攻击者无法获得 PA 验证预言机为显式前提。针对侧信道的防御与本文正交，可同时部署。

\textbf{物理攻击与固件层攻击。} 冷启动攻击、DMA 攻击、对固件或引导链的攻击不在范围内，它们违反第 2.6.3 节的信任基础。

\textbf{拒绝服务。} 攻击者主动触发大量随机化以制造性能下降，属于可用性攻击。本文在第 4.5.3 节给出各类负载下的触发频率，不将拒绝服务作为主要防御目标。

\textbf{关于内核符号表接口的说明。} \lstinline[breaklines=false]@/proc/kallsyms@ 与 \lstinline[breaklines=false]@dmesg@ 等接口会暴露内核符号地址，是内核中最常见的地址泄露途径。\textbf{本文的设计使这一途径不构成对不变式 I1 的威胁}：经第 3 章的编译改造之后，被随机化函数的符号绑定在其固定位置的跳板上，符号表中记录的是跳板地址而非函数在随机化区域内的实际地址。攻击者读取符号表只能得到跳板地址，而跳板位置本就是公开的，读到它不减少猜中随机化区域布局的难度。因此本节不把符号表接口列为排除项，而把它视为已被结构性覆盖的一类。

\section{安全不变式}

由第 2.6 节的威胁模型，本文提出三条安全不变式。三条不变式共同刻画了本文系统的防御目标，第 3、4、5 章的安全性分析分别说明本章机制如何维持它们。

\subsection{不变式的定义}

\textbf{I1，机密性。} 攻击者无法获知随机化区域内代码的当前地址。形式地说，对随机化区域内的任一函数 f，攻击者在任一时刻对 f 当前地址的最优猜测，其成功概率不超过在随机化区域可用地址空间上均匀猜测的概率。这条不变式针对向量 V1 与 V3：V1 被第 4 章的代码读取检测所覆盖，V3 被第 3 章的跳板机制在结构上消除，因为随机化之后内存中不再存在直接指向随机化区域内部的索引，所有索引都指向位置固定的跳板。

\textbf{I2，时效性。} 从任一随机化线索被泄露，到该线索失效，攻击者可完成的控制流转移次数有上界。I1 是一个理想目标，实践中无法保证攻击者一次也无法获得任何线索。I2 是 I1 的现实退化版本：承认线索可能泄露，但要求泄露到失效之间的窗口足够窄，使攻击者无法在窗口内凑齐可用的 gadget 集合。这条不变式针对向量 V2，由第 4 章的检测与触发机制维持，其可行性由第 3 章的低开销随机化机制保障。没有足够低的单次随机化代价，I2 所要求的响应速度就无法达到，这是第 3、4 两章相互依赖的形式表述。

\textbf{I3，完整性。} 任何跨区域的间接控制流转移，其目标必须属于该调用点的合法目标集合。这条不变式针对向量 V4。V4-A 由 PA 对区域内间接转移的验证约束；V4-B 由跳板处的分区验证约束。由第 5 章维持。

\subsection{不变式与章节、证据层级的对应}

\begin{table}[!htbp]
    \bicaption{\enspace 安全不变式的维持与论证方式}{\enspace Maintenance of the Security Invariants and the Form of Their Argument}
    \label{tab:2-5}
    \centering
    \scriptsize
    \setlength{\tabcolsep}{3pt}
    \renewcommand{\arraystretch}{1.25}
    \hyphenpenalty=10000\exhyphenpenalty=10000\emergencystretch=1em
    \begin{tabular}{>{\raggedright\arraybackslash}p{\dimexpr 0.0509\textwidth-2\tabcolsep\relax}>{\raggedright\arraybackslash}p{\dimexpr 0.0555\textwidth-2\tabcolsep\relax}>{\raggedright\arraybackslash}p{\dimexpr 0.1988\textwidth-2\tabcolsep\relax}>{\raggedright\arraybackslash}p{\dimexpr 0.1064\textwidth-2\tabcolsep\relax}>{\raggedright\arraybackslash}p{\dimexpr 0.5734\textwidth-2\tabcolsep\relax}}
        \hline
        不变式 & 针对向量 & 主要维持章节 & 依赖 & 论证的主要证据层级（见第 1.5 节）\\
        \hline
        I1 机密性 & V1, V3 & 第 4 章（检测）\newline 第 3 章（结构消除 V3） & — & L1：跳板是唯一出入口，可机械验证\newline L2：随机化熵与猜测成功概率\newline L3：检测覆盖率实测见第 4 章\\
        I2 时效性 & V2 & 第 4 章 & 依赖第 3 章的低开销 & L3：触发延迟分布，与随机化间隔上界对比\newline L4：未覆盖的试探途径显式列举\\
        I3 完整性 & V4-A, V4-B & 第 5 章 & 依赖第 3 章的跳板边界 & L1：跨区域转移必经跳板，可机械验证\newline L2：签名伪造概率（条件性）\newline L4：PA 验证预言机前提\newline L3：跨区域劫持的拦截见第 5 章\\
        \hline
    \end{tabular}
\end{table}

三条不变式中，只有“跳板是随机化区域的唯一出入口”这一条结构性质可以被完全验证，而 I1 与 I3 的论证都以它成立为前提。因此第 3 章给出的验证程序是全文安全论证的地基，其验证方法与扫描结果均在第 3 章给出。

\section{本章小结}

本章综述了与本文相关的研究工作，并给出了贯穿全文的威胁模型与安全不变式。

在研究现状方面，本章沿内存安全攻防的分阶段模型梳理了防御与攻击交替演进的脉络，指出内核随机化经历了一次性随机化、持续随机化、部分与触发式随机化三个阶段，每一阶段都在弥补前一阶段被攻破的假设，而当前所处的第三阶段以牺牲范围或牺牲函数位置的可变性换取了可行性，由此留下了 built-in 代码的保护盲区与非随机化区域的新攻击面。围绕持续随机化必须回答的三个核心问题，本章综述了已有解法：如何随机化一问对应代码索引更新与执行流的安全暂停两个方向，何时随机化一问对应触发策略，范围之外如何保护一问对应非随机化区域的处理。本章指出这些解法在 built-in 内核代码场景下各自失效的原因。本章还综述了本文两项关键技术手段的现状：只执行内存在内核中面临代码页多重映射与合法自读路径两个困难，而硬件辅助 CFI 的三类机制存在一个共同的精度上限，即硬件只能校验常数时间内可判定的谓词。

在威胁模型方面，本章定义了攻击者的四项能力、四类攻击向量、四项信任基础与四类明确的排除项，并说明了内核符号表接口为何不构成对本文防御目标的威胁。在此基础上导出三条安全不变式：机密性 I1、时效性 I2 与完整性 I3。三条不变式分别由后续三章维持，其相互依赖关系也在此明确：I2 的可行性依赖第 3 章提供的低开销随机化机制，I3 的论证依赖第 3 章建立的跳板边界。这一依赖关系是本文三项工作彼此耦合的形式表述，也是后续章节按依赖顺序编排的理由。

